\section{Domain Protocol Cards}
\label{sec:ieee-app-protocols}

The main paper compresses the active three domain rows into one shared narrative. This appendix
records the protocol-card view used during revision so that each row remains tied to one
problem statement, one safety object, and one closure boundary.

\begin{table*}[h]
\centering
\caption{Protocol-card summary for the active ORIUS domain rows.}
\label{tab:ieee-appendix-protocols}
\small
\begin{tabularx}{\textwidth}{@{}p{0.15\textwidth}p{0.22\textwidth}p{0.16\textwidth}p{0.18\textwidth}p{0.21\textwidth}@{}}
\toprule
\textbf{Domain} & \textbf{Canonical safety object} & \textbf{Repair mode} & \textbf{Fallback mode} & \textbf{Current closure note}\\
\midrule
Battery & State-of-charge and dispatch envelope & One-dimensional projection / clipping & Safe hold & Witness row with deepest runtime and theorem closure.\\
Autonomous vehicles & TTC and entry-barrier preservation & Bounded longitudinal repair & Full brake / bounded brake & Defended all-zip grouped nuPlan replay row under the current TTC entry-barrier contract.\\
Healthcare & Vital-alert and intervention envelope & Bounded intervention projection & Max alert & Defended bounded MIMIC-backed monitoring row; no live clinical deployment claim.\\
\bottomrule
\end{tabularx}
\end{table*}
